%% setting.tex — Problem setting, assumptions, definitions
%% Fixes Issues: #1 (no proofs), #2 (strong convexity), #3/#4 (gauge),
%%               #12 (regret definition), #13 (probability mode)
\section{Problem Setting and Assumptions}
\label{sec:setting}

%% ======================================================================
\subsection{Streaming Low-Rank Learning}
%% ======================================================================

We consider a learner receiving a stream of data $(x_t, y_t)_{t=1}^T$
drawn i.i.d.\ from a time-varying distribution $\mathcal{D}_t$. The
learner maintains a parameter $\theta_t \in \R^d$ and suffers per-round
loss $\ell_t(\theta) \coloneqq \ell(\theta; x_t, y_t)$.

\begin{definition}[Population loss]
\label{def:population-loss}
The \emph{population loss} at time $t$ is
\[
  L_t(\theta) \coloneqq \E_{(x,y) \sim \mathcal{D}_t}\bigl[\ell(\theta; x, y)\bigr].
\]
\end{definition}

\begin{definition}[Ground-truth optimum]
\label{def:optimum}
The time-varying optimum is $\theta^*_t \coloneqq \argmin_{\theta \in \R^d} L_t(\theta)$,
assumed to exist and be unique (ensured by strong convexity below).
\end{definition}

%% ======================================================================
\subsection{Two-Rate Structural Decomposition}
%% ======================================================================

\begin{assumption}[Low-rank decomposition]
\label{ass:decomposition}
There exist a slowly varying subspace $\mathcal{U}_t = \col(U_t)$ where
$U_t \in \R^{d \times r}$ has orthonormal columns ($U_t^\top U_t = I_r$,
so $U_t$ represents a point on the Grassmannian $\Gr(d,r)$), a fast-varying
coefficient vector $a_t \in \R^r$, and a residual $\rho_t \in \R^d$
such that
\[
  \theta^*_t = U_t\, a_t + \rho_t,
  \qquad
  \norm{\rho_t} \leq \rho_{\max},
  \qquad
  1 \leq r < d.
\]
The residual $\rho_t$ captures the approximation error of projecting
$\theta^*_t$ onto the rank-$r$ subspace.
\end{assumption}

%% ======================================================================
\subsection{Gauge-Invariant Subspace Distance}
%% ======================================================================

\begin{remark}[Gauge invariance --- fixes Issues \#3, \#4]
\label{rem:gauge}
The Frobenius distance $\norm{U_t - U_{t'}}_F$ is \emph{not}
gauge-invariant: for any orthogonal $R \in \mathcal{O}(r)$,
$U_t' = U_t R$ represents the \emph{same} subspace but can have
$\norm{U_t - U_t'}_F$ up to $\sqrt{2r}$.
We therefore use the \emph{principal-angle (geodesic) distance}
on $\Gr(d,r)$.
\end{remark}

\begin{definition}[Grassmann distance]
\label{def:grassmann-dist}
Let $U, U' \in \R^{d \times r}$ have orthonormal columns. Let
$\sigma_1 \geq \cdots \geq \sigma_r$ be the singular values of
$U^\top U'$. The \emph{principal angles} are
$\theta_i = \arccos(\min(\sigma_i, 1))$ for $i = 1, \ldots, r$.
The Grassmann distance is
\[
  \dG(\col(U),\, \col(U'))
  \coloneqq \norm{\sin \Theta}_F
  = \Bigl(\sum_{i=1}^r \sin^2\theta_i\Bigr)^{1/2}.
\]
Equivalently, $\dG^2 = r - \norm{U^\top U'}_F^2$ when both have
orthonormal columns.
\end{definition}

\begin{lemma}[Grassmann distance properties]
\label{lem:grassmann-props}
$\dG$ is a metric on $\Gr(d,r)$ satisfying:
\begin{enumerate}[nosep]
  \item $0 \leq \dG \leq \sqrt{r}$;
  \item $\dG(\col(U), \col(UR)) = 0$ for all $R \in \mathcal{O}(r)$
        \textup{(gauge invariance)};
  \item $\dG \leq \norm{U - U'}_F \leq \sqrt{2}\, \dG$ when both
        $U, U'$ are orthonormal \textup{(Frobenius sandwich)}.
\end{enumerate}
\end{lemma}

\begin{proof}
(1) follows from $\sin^2\theta_i \in [0,1]$.
(2): $U$ and $UR$ have the same column space, so
$\sigma_i(U^\top(UR)) = \sigma_i(R) = 1$ for all $i$, giving
$\theta_i = 0$.
(3): Using $\norm{U - U'}_F^2 = 2r - 2\tr(U^\top U')$ and the
relation between trace and singular values:
$\tr(U^\top U') \leq \sum_i \sigma_i = \sum_i \cos\theta_i$,
so $\norm{U-U'}_F^2 = 2\sum_i(1-\cos\theta_i) \leq
2\sum_i \sin^2\theta_i / (1+\cos\theta_i) \cdot 2 \leq
2\sum_i \sin^2\theta_i = 2\dG^2$ using $1-\cos\theta \leq
\sin^2\theta$ for $\theta \in [0,\pi/2]$.
The lower bound: $\norm{U-U'}_F^2 \geq 2\sum_i(1-\cos\theta_i)
\geq \sum_i \sin^2\theta_i = \dG^2$ using $1-\cos\theta \geq
\sin^2\theta/2$.
\end{proof}

%% ======================================================================
\subsection{Drift Variation Budgets}
%% ======================================================================

\begin{definition}[Slow variation]
\label{def:slow-variation}
The \emph{slow (subspace) variation} over horizon $T$ is
\[
  V_s(T) \coloneqq \sum_{t=1}^{T-1} \dG\bigl(\col(U_t),\, \col(U_{t+1})\bigr).
\]
\end{definition}

\begin{definition}[Fast variation --- jump process]
\label{def:fast-variation}
The fast (coefficient) variation is modeled as a \emph{piecewise-constant
jump process}: there exist times $0 = \tau_0 < \tau_1 < \cdots < \tau_{K_T} \leq T$
such that $a_t = a_{\tau_k}$ for $t \in [\tau_k, \tau_{k+1})$. We write
$K_T$ for the number of jumps and define
\[
  V_f(T) \coloneqq \sum_{k=1}^{K_T} \norm{a_{\tau_k} - a_{\tau_{k-1}}}.
\]
\end{definition}

\begin{assumption}[Bounded variation]
\label{ass:bounded-variation}
$V_s(T)$ and $V_f(T)$ are finite for all $T$. We also assume
minimum jump size $\Delta_0 > 0$:
\[
  \norm{a_{\tau_k} - a_{\tau_{k-1}}} \geq \Delta_0
  \quad \text{for all } k = 1, \ldots, K_T.
\]
\end{assumption}

\begin{assumption}[Bounded coefficients]
\label{ass:bounded-coeff}
$\norm{a_t} \leq A_{\max}$ for all $t$. This ensures the jump magnitudes
are bounded: $\norm{a_{\tau_k} - a_{\tau_{k-1}}} \leq 2A_{\max}$.
\end{assumption}

%% ======================================================================
\subsection{Loss Assumptions (on Population Loss)}
%% ======================================================================

The following assumptions apply to the \emph{population loss} $L_t$,
\textbf{not} the per-sample loss $\ell_t$. This is a critical distinction
(cf.\ Issue~\#2): the per-sample squared loss
$\ell_t(\theta) = \frac{1}{2}(\theta^\top x_t - y_t)^2$ has rank-1
Hessian $x_t x_t^\top$ for $d > 1$, which is \emph{not} strongly convex.
But the population loss $L_t(\theta) = \frac{1}{2}\E[(\theta^\top x - y)^2]$
has Hessian $\Sigma_t = \E[x_t x_t^\top]$, which is $\mu$-strongly convex
when $\lambda_{\min}(\Sigma_t) \geq \mu > 0$.

\begin{assumption}[Strong convexity of population loss]
\label{ass:strong-convexity}
For all $t \in [T]$, the population loss $L_t$ is $\mu$-strongly convex:
\[
  L_t(\theta') \geq L_t(\theta) + \inner{\nabla L_t(\theta)}{\theta'-\theta}
  + \frac{\mu}{2}\norm{\theta'-\theta}^2,
  \qquad \forall\, \theta, \theta' \in \R^d,
\]
where $\mu = \inf_{t \in [T]} \lambda_{\min}\bigl(\nabla^2 L_t(\theta^*_t)\bigr) > 0$.
\end{assumption}

\begin{assumption}[$L$-smoothness of population loss]
\label{ass:smoothness}
For all $t \in [T]$, $L_t$ is $L$-smooth:
\[
  \norm{\nabla L_t(\theta) - \nabla L_t(\theta')} \leq L\norm{\theta - \theta'},
  \qquad \forall\, \theta, \theta' \in \R^d.
\]
The condition number is $\kappa \coloneqq L/\mu \geq 1$.
\end{assumption}

\begin{assumption}[Bounded noise]
\label{ass:noise}
The stochastic gradient $g_t(\theta) = \nabla \ell_t(\theta)$ is an
unbiased estimator of $\nabla L_t(\theta)$ with bounded variance:
\[
  \E\bigl[\norm{g_t(\theta) - \nabla L_t(\theta)}^2 \bigm| \theta\bigr]
  \leq \sigma^2, \qquad \forall\, \theta \in \R^d.
\]
\end{assumption}

%% ======================================================================
\subsection{Regret Definition}
%% ======================================================================

\begin{definition}[Expected dynamic regret --- fixes Issue~\#12]
\label{def:regret}
The \emph{expected dynamic regret} of an algorithm producing
$\theta_1, \ldots, \theta_T$ (where $\theta_t$ is $\mathcal{F}_{t-1}$-measurable,
$\mathcal{F}_t = \sigma(x_1,y_1,\ldots,x_t,y_t)$) is
\[
  \reg_T \coloneqq \E\Biggl[\sum_{t=1}^T
    \bigl(L_t(\theta_t) - L_t(\theta^*_t)\bigr)\Biggr],
\]
where the expectation is over the randomness in $(x_t, y_t)_{t=1}^T$
and any internal randomness of the algorithm.
\end{definition}

\begin{remark}[Probability mode --- fixes Issue~\#13]
\label{rem:prob-mode}
All bounds in this supplement are stated \emph{in expectation}.
By Markov's inequality, an expected regret bound of $\E[\reg_T] \leq R$
implies a high-probability bound $\Pr[\reg_T > R/\delta] \leq \delta$.
Sharper high-probability bounds (e.g., via Azuma--Hoeffding for bounded
differences) are possible but beyond the scope of these proofs.
\end{remark}

%% ======================================================================
\subsection{Applicability to LLM Fine-Tuning}
%% ======================================================================

\begin{remark}[Theory--practice scope]
\label{rem:scope}
\Cref{ass:strong-convexity,ass:smoothness} hold for quadratic / linear
regression with well-conditioned covariance. They do \emph{not} hold
for LLM next-token prediction loss, which is non-convex in the full
parameter space and only approximately locally convex in the LoRA subspace
near a pre-trained minimum.

The theorems below therefore provide \emph{motivating theory} for the
dual-timescale principle. The Mirror-LoRA algorithm applied to LLMs is
a \emph{theory-inspired heuristic}, not a formally justified instantiation
of the theoretical framework. The empirical evaluation (Section~5 of
the main paper) demonstrates that the dual-timescale principle yields
strong performance on real LLMs despite the gap between theory and practice.
\end{remark}

%% ======================================================================
\subsection{Assumption Summary}
%% ======================================================================

\begin{center}
\begin{tabular}{clcc}
\toprule
\textbf{ID} & \textbf{Assumption} & \textbf{Used by} & \textbf{Status} \\
\midrule
A1 & $L_t$ is $\mu$-strongly convex (\Cref{ass:strong-convexity})
   & Thm~1,2 & Holds for linear-Gaussian \\
A2 & $L_t$ is $L$-smooth (\Cref{ass:smoothness})
   & Thm~1,2 & Holds for linear-Gaussian \\
A3 & Low-rank decomposition (\Cref{ass:decomposition})
   & Thm~1,2 & By construction (synthetic) \\
A4 & Bounded variation (\Cref{ass:bounded-variation})
   & Thm~1,2 & Finite horizon \\
A5 & Bounded coefficients (\Cref{ass:bounded-coeff})
   & Thm~1 & Replacement jumps \\
A6 & Bounded noise (\Cref{ass:noise})
   & Thm~2 & Bounded loss \\
\bottomrule
\end{tabular}
\end{center}
